\section{Problem 3: Gear Train Design}

\subsection{Problem description}
The objective of this exercise is to choose a gear train design that achieves a specific transmission ratio.
We would define the gear train with four gears with the numbers of teeth $N_1, N_2, N_3, N_4$.
The target transmission ratio is:
\begin{equation}
R = \frac{N_2 \cdot N_4}{N_1 \cdot N_3}
\end{equation}
(Note: Standard formulation is usually $\frac{product\ driven}{product\ drivers}$ or similar. The equation in text is $R=(N2 N4)/(N1 N3)$ or $f(N) = (N1 N3)/(N2 N4)$. I will use the cost function definition for consistency: $Target \approx 1/6.931 \approx 0.144$).

Here, the target transmission ratio will be randomly chosen: $R_{target} = 1/6.931 \approx 0.1442793$.
And each gear must have an integer number of teeth between 12 and 60.
The problem is difficult because the search space is discrete and nonlinear, meaning classical optimization cannot be used directly.
Therefore, two heuristic optimization algorithms will be tested:
\begin{itemize}
    \item Particle Swarm Optimization (PSO)
    \item Simulated Annealing (SA)
\end{itemize}
The goal is to compare their performance and their behavior on this discrete optimization problem.

\subsection{Cost function}
The cost function measures the error between the achieved ratio and the target ratio:
\begin{equation}
f(N_1,N_2,N_3,N_4) = \left| \frac{N_1 \times N_3}{N_2 \times N_4} - R_{target} \right|
\end{equation}
A good solution is one where this error approaches zero.

\subsection{Sources}
Course material on PSO and Simulated Annealing.

\subsection{Methodology}

\textbf{Particle Swarm Optimization:}
It simulates a swarm of particles exploring the search space. Each particle represents a potential gear combination $[N_1,N_2,N_3,N_4]$.
Particles move according to:
\begin{itemize}
    \item Their own best-found solution
    \item The best solution found by the swarm
\end{itemize}
Positions are rounded to integers to satisfy the gear teeth constraint as we can’t make gears with an odd number of teeth.
We store the best error at each iteration to generate a convergence curve like we did for the SA.

\textbf{Simulated Annealing:}
It explores the space by randomly perturbing the current solution. It accepts worse solutions with a probability depending on analogy to a temperature parameter:
\begin{itemize}
    \item With high temperature goes a wide exploration
    \item Low temperature goes a convergence around a good solution
\end{itemize}
The algorithm stops when the temperature becomes very low. We record both the current energy and the best energy to produce a convergence curve.

\subsection{Settings}

\textbf{PSO}
\begin{itemize}
    \item Population size: 40
    \item Iterations: 300
    \item $w=0.7, C_1=1.5, C_2=1.5$
    \item Integer rounding of particle positions
    \item Recorded global best error per iteration (for convergence curve)
\end{itemize}

\textbf{SA}
\begin{itemize}
    \item Initial temperature: 1000
    \item Cooling rate: 0.90
    \item Minimum temperature: 0.1
    \item Random seed: 42
    \item Recorded: Current energy / Best energy (for convergence curve)
\end{itemize}

\subsection{Results}

\subsubsection{Results analysis}
The optimal solution found by the Particle Swarm Optimization is:

\textbf{PSO Results:}
PSO finds an extremely accurate gear ratio, with an error on the order of $10^{-5}$, meaning it almost exactly matches the target ratio.

\begin{figure}[h!]
    \centering
    \includegraphics[width=0.8\textwidth]{Gear_Train_PSO_Convergence.png}
    \caption{PSO Convergence curve}
\end{figure}

The graph shows that PSO converges extremely quickly within $\approx 10$ iterations and then stabilizes at a very small error value.

And the result found by the Simulated Annealing (Solid Annealing) is:

\textbf{SA Results:}
SA clearly fails on this problem. The final ratio is far from the target, with a very large error.

\begin{figure}[h!]
    \centering
    \includegraphics[width=0.8\textwidth]{Gear_Train_SA_Convergence.png}
    \caption{SA Convergence Curve}
\end{figure}

The plot shows:
\begin{itemize}
    \item Highly unstable exploration, current energy jumps up and down
    \item Best energy stays flat because the algorithm never discovers a good solution
    \item Cooling happens too fast to explore thoroughly
\end{itemize}

\subsubsection{Explanations}

\textbf{Why PSO succeeds}
\begin{itemize}
    \item PSO is particularly effective in structured continuous search spaces, even with integer rounding.
    \item The particles rapidly converge to the correct ratio because the cost function is smooth.
    \item Global and personal best mechanisms guide the search efficiently.
\end{itemize}

\textbf{Why SA fails}
\begin{itemize}
    \item The gear train problem is very sensitive: small changes in $N_i$ produce large jumps in the ratio.
    \item SA explores randomly and relies on thermal acceptance, which:
    \begin{itemize}
        \item Moves too chaotically
        \item Does not exploit promising regions enough
        \item Requires much slower cooling to perform well
    \end{itemize}
    \item SA found a physically valid gear set but not even close to the desired ratio.
\end{itemize}

\subsubsection{Conclusion}
PSO is clearly the superior method for this task. It finds a nearly perfect solution extremely fast, while SA remains stuck far away from any optimal region.
For a future study, SA could be improved by:
\begin{itemize}
    \item Reducing the cooling rate from 0.99 to 0.90 for example
    \item Using a smarter neighbour generator
    \item Introducing integer-specific heuristics
\end{itemize}
But in its current form, SA is not suitable for this optimisation problem.
